\section*{The fixed-axis arithmetic jet in the full orthogonal correspondence}

\subsection*{The precise conversion between the retained pair and triple}

Use the complete Hilbert spaces, domains, reality conditions and operators
of the preceding solenoidal and swirl appendices. In particular
\(P_\nu=K_\nu K_\nu^*\), \(\Pi\mathcal R=I\),
\(\mathcal R^*=\Pi\), and \(\Pi K_\nu=0\).
Let \(w_{\rm axis}\) denote exactly the \(w_\nu=u_\nu-\mathcal R B_\nu\)
in the preceding axis appendix. It has the typed domain
\[
 w_{\rm axis}\in V_\sigma\cap\ker\Pi,
 \qquad B_\nu\in\mathcal B.
\]
The complete conversion from that pair to the solenoidal triple is
\begin{align*}
 (B_\nu,w_{\rm axis})&\longmapsto
   (h_{\rm sol},B_\nu,v),\\
 h_{\rm sol}&=K_\nu^*w_{\rm axis},
 &v&=(I-P_\nu)w_{\rm axis},\\
 (h_{\rm sol},B_\nu,v)&\longmapsto
   (B_\nu,K_\nu h_{\rm sol}+v),
 &v&\in V_\sigma\cap\ker K_\nu^*\cap\ker\Pi.
\end{align*}
Indeed \(K_\nu^*\mathcal R=(\Pi K_\nu)^*=0\), so
\(K_\nu^*w_{\rm axis}=K_\nu^*u_\nu\).
Also \(\Pi P_\nu=0\), and hence \(v\) is in the displayed kernel
intersection; \(K_\nu^*(I-P_\nu)=0\) supplies its other kernel condition.
Adding \(P_\nu w_{\rm axis}\) and \((I-P_\nu)w_{\rm axis}\)
proves the first inverse identity. Conversely, applying \(K_\nu^*\)
and \(I-P_\nu\) to \(K_\nu h_{\rm sol}+v\) returns
\(h_{\rm sol}\) and \(v\), respectively. Thus both maps are onto their
displayed targets and are mutual inverses. The full reconstruction and
energy are still exactly
\[
 u_\nu=K_\nu h_{\rm sol}+\mathcal R B_\nu+v,
 \qquad
 \|u_\nu\|_V^2=
 \|h_{\rm sol}\|_H^2+\|B_\nu\|_{\mathcal B}^2+\|v\|_V^2.
\]
The original solenoidal pair remainder is instead
\(w_{\rm sol}=(I-P_\nu)u_\nu=\mathcal R B_\nu+v\).
Its exact conversion is \(B_\nu=\Pi w_{\rm sol}\) and
\(v=(I-\mathcal R\Pi)w_{\rm sol}\).
This specifies the relation between both remainder coordinates without
discarding either one. Applying the arithmetic operator to the same
\(B_\nu\) therefore composes the two correspondences literally, with
the original \(\sigma=r^2/2\), physical \(z,t,\nu\), all nine
nonlinear ordered pairs, and the previously displayed pressure and force.

\subsection*{An exact on-axis vorticity functional and its arithmetic image}

The endpoint functional is defined on smooth physical fields, not on
arbitrary \(L^2\) equivalence classes. Fix \(z\), orient the
\((x_1,x_2)\)-disk by \(+e_3\), and orient its circle by increasing
\(\theta\). For a smooth vector field \(u\), write
\(\omega_3=\partial_{x_1}u_2-\partial_{x_2}u_1\).
The original angular-momentum projection has the exact circulation form
\begin{align*}
 (\Pi u)(r^2/2,z)
 &=\frac1{2\pi}\oint_{|x_\perp|=r}u\cdot d\ell\\
 &=\frac1{2\pi}\int_{|x_\perp|\le r}\omega_3(x_\perp,z)\,dx_1dx_2
 =\int_0^r \rho\,\langle\omega_3(\rho,\theta,z)\rangle\,d\rho.
\end{align*}
The first equality uses \(d\ell=r e_\theta\,d\theta\); the second is
Stokes' formula with the stated orientation. For \(r>0\), differentiating
and using \(d\sigma/dr=r\) gives
\[
 \partial_\sigma(\Pi u)(r^2/2,z)
       =\langle\omega_3(r,\theta,z)\rangle.
\]
Continuity at the axis gives the endpoint value: the difference from
\(\omega_3(0,0,z)\) is bounded in absolute value by
\(\sup_{|x_\perp|\le r}|\omega_3(x_\perp,z)-\omega_3(0,0,z)|\),
which tends to zero. The integral formula also proves the one-sided
derivative at \(\sigma=0\) directly. Consequently the typed functional
\[
 \mathcal J_z(u):=\partial_\sigma(\Pi u)(0,z)
       =\omega_3(0,0,z)
\]
annihilates every smooth field in \(\ker\Pi\). In particular the
smooth preterminal components satisfy
\[
 \mathcal J_z(K_\nu h_{\rm sol})=0,\qquad
 \mathcal J_z(v)=0,\qquad
 \mathcal J_z(u_\nu)=\mathcal J_z(\mathcal R B_\nu).
\]
Smoothness here follows from the full Fourier-domain preservation proved
earlier: the actual source lies in every spatial Sobolev domain at each
preterminal time, and each of the three orthogonal projections preserves
those domains. Sobolev embedding at successively larger orders gives
smooth representatives for all three components.

For the compact smooth source field, the complete endpoint summation
proof above now gives the exact linear-functional identity
\[
 \left.\partial_x\mathcal A_h(\Pi u_\nu)(x,z,t)\right|_{x=0}
    =\frac{2^{h+1}-1}{16}\,\mathcal J_z(u_\nu)
    =\frac{2^{h+1}-1}{16}\,\mathcal J_z(\mathcal R B_\nu).
\]
Thus the arithmetic endpoint functional acts on the explicit orthogonal
swirl component. For the actual source at \(z=0,t=1-\tau\), the proved
axis identity supplies, on its original interval \(0<\tau<\tau_1\),
\begin{align*}
 \omega_{3,\nu}(0,0,0,1-\tau)&=\frac2C\tau^{-1-h},\\
 \partial_x\mathcal A_hB_\nu(0,0,1-\tau)
 &=\operatorname{Res}_{s=-1}\mathfrak B_\nu(s,0,1-\tau)
  =\frac{2^{h+1}-1}{8C}\tau^{-1-h}.
\end{align*}
The source amplitude \(C>0\) and exponent \(h\) retain their original
values; the complete viscosity map proves why no additional \(\nu\)
factor remains in this fixed-axis derivative. The exact inner-circle
velocity estimate, which has its own \(\sqrt\nu\) and original remainder,
remains the different evaluation already proved in the preceding appendix.

The fixed-axis result and every endpoint derivative were received as a
complete source proof, preserved in \path{sources/ns_axis_jet/parent_29l.tex}.
The accompanying review records its precise reviewed hash; the provenance
file records the frozen source hash and reader adaptations separately.
The complete Euler--Maclaurin argument needed for the adaptation is
included above, together with a direct derivation of \(\zeta(-1)\).
The present pair/triple conversion and circulation calculation supply the
exact composition with the full solenoidal state. The entire preceding
solenoidal-swirl extension also received a complete independent incoming
owner review with no mathematical correction; its final physical/Fourier
operator notation clarification was checked locally.
